\documentclass[11pt]{article}
\usepackage[T1]{fontenc}
\usepackage{amsmath,amssymb,geometry}
\geometry{margin=27mm}
\title{EPPA III: fragile combinatorial arguments (first audit)}
\date{16 October 2026}
\begin{document}\maketitle
\paragraph{Scope.}
Source: \texttt{main-new(20261016-172438).tex}. We accept external
Gardiner, rank-three, and local-recognition classifications as inputs.
These are paper-level checks, not new Lean certificates.
The original manuscript is unchanged.

\paragraph{1. Lemma 7.6: lift of partial automorphisms.}
The displayed map $\widehat\alpha(i,d)=(i,\alpha(d))$ need not
take its values in the selected $\widehat G$.
For example, in a two-clique extension of $H=P_4$, select only
$(1,v),(2,w)$ over adjacent $v,w$. The transposition of
$v,w$ sends $(1,v)$ to the unselected $(1,w)$.
The lemma can be repaired without changing its statement:
choose a vertex $s(v)\in\widehat G$ above each occupied
$v\in\pi(\widehat G)$. Lift any partial automorphism
$\alpha$ by $s(v)\mapsto s(\alpha v)$. This is a partial
automorphism of $\widehat G$ because relations between
distinct fibres depend only on their projections.
An EPPA extension preserves the true-twin fibres of
$\widehat H$; as $H$ has no twins these classes are
exactly $[l]\times\{v\}$ (singletons if $l=1$).
The induced automorphism of $H$ extends $\alpha$.
In the preceding paragraph, $\pi(S)$ is an induced subgraph
of $H$, not of $G$.

\paragraph{2. Example 8.4: arithmetic.}
The local graph $T(24)$ has order $\binom{24}{2}=276$,
not $138$. Excluding $|G_1|\ge134$ gives $|G_1|\le133$,
not $<133$. Then
\[
 |G_2|\ge1024-1-133=890>(2048-1-276)/2=|H_2|/2.
\]
For $|G_1|\ge134>24$, Remark 4.1 also supplies the
triangle-free line-graph hypothesis of Theorem 4.7.
These repairs preserve the intended contradiction.

\paragraph{3. Lemma 7.9: strongly regular generalised covers.}
The conclusion is supported, but the original proof contains
false identifications. For $s=1$, $X=\overline M$ gives
$\overline{tK_2}$ for all $t$, not only $t\le2$.
For $t=1$, $X=Y=1$ gives $K_{s,s}$, not $sK_2$.
Here is a complete short replacement proof.

Assume $s,t\ge2$. Let $q$ be the number of common
neighbours in $D_i$ of two distinct vertices of $C_i$.
For $X=0,1,M,\overline M$, the values are $0,t,0,t-2$.
If $Y=0$, two vertices in distinct $C_i$ have no
common neighbours, so $\mu=0$. If $X$ is $M$ or
$\overline M$, a missing edge inside $C_i\cup D_i$
has a common neighbour, contradicting $\mu=0$.
Hence $X=0$ or $1$, giving $2sK_t$ or $sK_{2t}$.

If $Y=1$, two adjacent vertices of $C_i$ have
$t-2+q+(s-1)t$ common neighbours, whereas an edge
between $C_i$ and $D_j$, $j\ne i$, has $2(t-1)$.
Therefore
\[
 t-2+q+(s-1)t=2(t-1),\quad q=(2-s)t.
\]
Hence $s=2$ and $q=0$. The only possibilities are
$X=0$, $X=M$, or $X=\overline M$ for $t=2$.
The latter two have a diagonal matching edge
with zero common neighbours, contradicting the
$2t-2>0$ common neighbours of an off-diagonal edge.
Thus $X=0$, giving $2K_{2t}$.

For $s=1$, the candidates are $2K_t$, $K_{2t}$,
$\overline{tK_2}$, and two cliques joined by a matching.
In the last case strong regularity forces $t\le2$.
For $t=1$, the host is bipartite with empty,
complete, matching, or complementary matching cross edges.
The first three are homogeneous; the last is strongly
regular only when $s\le2$, and is then also homogeneous.
Complementing yields the complementary case.

\paragraph{4. Example 8.3: EPPA orbit justification.}
A stabilizer acting on the set of non-neighbours does not
alone justify division by the cardinality of the set.
EPPA gives the missing transitivity: vertices of $G$
nonadjacent to all chosen independent selected vertices
are conjugate by a pointwise-stabilizing automorphism.
Thus the successive orbit sizes are at least $76,50,24$.
After choosing $n_1,n_2$, the stabilizer has order
at most $17280/(76\cdot50)<5$, whereas the next orbit
has at least 24 elements, a contradiction.

\paragraph{5. Lemma 5.2: partial two-graph switching.}
The original induction requests a previously selected
pair $(u,fu)$ even when starting the first nonperiodic
path. An alternative parity proof uses only the cited
Mallows--Sloane theorem.
For a base adjacency indicator $e(x,y)\in\mathbb F_2$,
put $\delta(x,y)=e(x,y)+e(fx,fy)$ on the domain $D$.
Preservation of the two-graph gives zero sum around
each triple, hence $\delta(x,y)=h(x)+h(y)$ for some
$h:D\to\mathbb F_2$.
On the cyclic points of $f$, Mallows--Sloane yields a
switching $s$ and a constant $c$ with
$s(x)+s(fx)=h(x)+c$. Choose $c$ arbitrarily if
there are no cycles. Extend $s$ over each remaining
directed path by the same equation, choosing its
first value freely. Then the switched adjacency
satisfies $e_s(x,y)=e_s(fx,fy)$ for all $x,y\in D$.
An exhaustive finite audit on five labelled vertices
tests 98,944 partial bijections across all 64 two-graphs;
all 44,104 partial two-graph automorphisms admit
a compatible switching. This is not Lean verification.

\paragraph{Follow-up.}
Fix Lemma 7.6, Example 8.4, Lemma 7.9, and the
EPPA-orbit wording in Example 8.3 first.
The remaining imprimitive classification obligations
are tracked separately.
\end{document}
